\section{Introduction}

Online bipartite matching is the canonical model of irrevocable
sequential allocation: offline resources are known in advance, requests
arrive one at a time, and each request must be matched to an available
neighbor --- or dropped --- immediately and irrevocably, before the rest
of the input is seen. It underlies ad allocation, ride-hailing dispatch,
organ exchange, and request routing in serving systems. Classical online
algorithms --- Greedy, KVV Ranking \cite{kvv1990ranking}, and the
stochastic-matching algorithms of Feldman et al.
\cite{feldman2009online} and Jaillet--Lu \cite{jailletlu2014online} ---
are analyzed through the worst-case competitive ratio.

A now-large body of work augments these algorithms with a
\emph{prediction} about the input and asks for two guarantees at once:
\textbf{consistency} (near-optimal performance when the prediction is
good) and \textbf{robustness} (no worse than an advice-free baseline
when the prediction is arbitrarily bad)
\cite{lykouris2018caching,weizhang2020tradeoffs}. For online matching
specifically, two strands have emerged. MinPredictedDegree (MPD)
\cite{aci2022mpd} matches each arrival to its available neighbor of
minimum \emph{predicted} degree, protecting scarce resources; it is
robust by construction (a useless predictor reduces it to Ranking). A
second strand --- the \emph{test-and-fallback} algorithms of Choo et
al.~\cite{choo2024imperfect} and Burathep--Erlebach--Moses
\cite{bem2026testmatch} --- is explicitly adaptive: it Mimics a
type-histogram prediction on a sublinear prefix of arrivals, tests
whether the observed arrivals match the prediction, and then either
continues to follow it or falls back to Ranking.

These algorithms have been studied \emph{in isolation} --- each on its
own input family, against its own error model, and largely through
worst-case theory. As a result, three basic questions have no empirical
answer. How do the algorithms actually compare, head to head, under one
prediction-error model? How much does a prediction help, and at what
cost, on realistic inputs? And --- most importantly --- \emph{why} does
the practical experience with these algorithms feel so different from
their worst-case promise? This paper answers all three, with a unified
experimental study and a theorem that explains what the study finds.

\textbf{A unified experimental study.} We place the advice-free
baselines, the MPD family (and its Feldman/Jaillet--Lu augmentations),
and the test-and-fallback algorithms on a single harness: common graph
families, a common structured prediction-error model, a common optimum
computed by Hopcroft--Karp, and 95\% confidence intervals throughout. We
also port the blind-follow-with-switching \emph{combiner} of Chłędowski
et al.~\cite{chledowski2021caching} --- the ``cheap worst-case
insurance'' from the caching literature --- and benchmark it (we do not
claim it as new). The study runs across synthetic graphs spanning the
difficulty range, six real-world graphs from the Network Repository, and
real request traces.

\textbf{The empirical wall.} Across every setting we observe the same
phenomenon, which we state as the paper's organizing thesis: \emph{on
average-case inputs the value of predictions is robustness insurance,
not a performance lever.} Concretely (Section 3): the advice-free
Ranking is already within a percent or two of the optimum, so the
consistency upside of even perfect advice is small; unguarded
prediction-following (blind Mimic, or MPD under adversarial predictions)
crashes \emph{below} the advice-free floor; and the entire practical
value of the sophisticated algorithms is that they refuse to crash. Two
distinct mechanisms deliver this --- the structural robustness of the
MPD-augmentations and the adaptive robustness of test-and-fallback ---
with different consistency/robustness trade-offs. We confirm the wall on
the six real graphs and on real traces, where a cheap, non-ML historical
predictor captures a meaningful fraction of the (small) oracle gap while
never dropping below the baseline (Section 6). We also report a negative
result: because MPD depends on its predictor only through the induced
\emph{order}, one might train the predictor with an order-aware loss
instead of regression, but on realistic features the two coincide ---
reinforcing that the predictor, once order-faithful, is not the
bottleneck.

\textbf{Why the wall is necessary --- and what it costs.} The empirical
wall is not an artifact of a particular generator or algorithm; it has a
price tag. Our main theoretical contribution (Section 7) is a two-sided
\textbf{budget--stakes law} for the test-and-fallback class. Informally:

\begin{quote}
Deciding whether to follow the advice costs a prefix of length
\(k^* = \tilde\Theta(\sigma^2/g^2)\), where \(g\) is the stakes (what
good advice gains over the baseline) and \(\sigma^2\), the arrival mass
on contested resources, equals exactly twice the baseline slack. Below
\(k^*\), \emph{no} decision rule is simultaneously consistent and
robust; at \(k^*\), an explicit one-line rule is.
\end{quote}

The two bounds meet up to logarithms whenever the stakes are carried, at
comparable signal levels, by a constant fraction of the contested mass
--- the condition under which ``\(k^*\)'' names a single quantity; §7.8
records the low-signal regime in which the two sides can still separate,
which we leave open. The lower half is information-theoretic --- it
binds \emph{any} measurable rule on the prefix, not merely the
empirical-distance threshold used in practice --- via a master
consistency/robustness inequality and a Hellinger computation. The upper
half is a \textbf{directional test}: classify each prefix arrival as
agreeing or disagreeing with the advice's local prediction, and follow
iff agreements win --- a rule we also implement, calibrate, and evaluate
head-to-head on the benchmark (§5.4). Its analysis rests on a
\emph{payoff identity} --- on the hard family, the advantage of
following is exactly half the mean of a per-sample observable --- which
also refutes the tempting stronger conjecture (and an earlier version of
our own claim) that tolerant-testing hardness
\cite{canonne2022tolerance} makes the decision impossible for every
sublinear rule: the advice's \emph{distance} to the truth is indeed hard
to test, and the field's empirical-\(\ell_1\) thresholds are provably
blind at sublinear prefixes, but the decision-relevant \emph{payoff} is
not. The wall then re-emerges exactly where the experiments live: stakes
are capped by the baseline slack
(\(g \le 2\varepsilon_{\max}(1-\rho_{\mathrm{base}})\), with
\(\varepsilon_{\max}\) the largest per-cell signal), so on
strong-baseline instances the budget \(k^*\) exceeds the horizon itself
--- every upside below \(\Theta(\sqrt{(1-\rho_{\mathrm{base}})/n})\) is
uncapturable at \emph{any} feasible prefix, and the upsides measured in
Sections 3--6 sit in that range.

\textbf{Contributions.} - \textbf{(C1) The first unified benchmark} of
learning-augmented online-matching algorithms --- advice-free baselines,
the MPD family, and the test-and-fallback schemes --- on one harness
with a common error model, a common optimum, and confidence intervals
(Section 3). - \textbf{(C2) The robustness-insurance characterization.}
Unguarded followers crash below the advice-free floor; two mechanisms
(structural and adaptive) restore safety with distinct trade-offs; the
consistency upside is small on average-case inputs, so the value is
downside protection. Validated on synthetic graphs, six real graphs, and
real traces (Sections 3, 6). - \textbf{(C3) An empirical engagement with
the order-error theory.} MPD's loss is governed by a Kendall-\(\tau\)
order error onto which several error models collapse; we characterize
the tightness and saturation of the known \(n{-}\mathrm{LIS}\) bound
{[}ACI22, App. D{]} rather than re-deriving order-dependence (Section
4). - \textbf{(C4) The first empirical study of test-and-fallback},
including a threshold-calibration pathology (a more accurate test can
make a worse decision), its recalibration and resolution limit, a
\textbf{constant-free payoff-testing rule} that avoids both threshold
pathologies --- its misjudgement \emph{falls} with prefix size exactly
where the thresholds' \emph{rises} --- and a benchmark of the dynamic
combiner exhibiting an irrevocability penalty that explains why matching
needs \emph{test-then-commit} (Section 5). - \textbf{(C5) The
budget--stakes law.} A two-sided law for test-and-fallback --- the two
sides meeting up to logarithms under a stated signal condition (§7.5,
§7.8): the follow/fallback decision costs a prefix of
\(\tilde\Theta(\sigma^2/g^2)\) (baseline slack over the square of the
stakes) for \emph{any} rule (lower bound), and an explicit directional
test achieves it (upper bound). A payoff identity separates cheap
payoff-testing from provably hard distance-testing --- refuting, and
reporting honestly, the natural tolerant-testing impossibility --- and
the corollary quantifies the wall: on strong-baseline instances, upsides
below \(\Theta(\sqrt{(1-\rho_{\mathrm{base}})/n})\) are uncapturable at
any prefix length (Section 7).

The experiments and the law are two views of one fact. The experiments
discover a wall --- predictions buy insurance, not performance; the law
prices the wall --- verifying advice costs the inverse square of its
stakes, and average-case stakes cannot cover the bill. We adopt an
AI-inference serving instantiation as a running application case study
(Section 8) rather than a novelty claim, and we are careful throughout
to credit the prior results our findings build on and to state the scope
of each claim.
